\chapter{Analisi}
\label{chap:analysis}

\section{Obbiettivi}
\label{sec:objectives}

Definizione chiara degli obbiettivi ed il focus di questo progetto:
\begin{itemize}
	\item Studio di Zenoh: il suo funzionamento, il modello peer
	      to peer (scouting, router, etc\dots) e come può essere utilizzato.
	\item Setup del techstack: embdedded rust, il framework esp-idf, integrazione
	      di librerie esterne e creazione del wrapper.
	      \sidenote{come dovrei usare inglesismi come techstack?}
	\item Creazione della network: comunicazione, gestione della perdita di nodi,
	      integrazione con yaair.
\end{itemize}

\section{Requisiti funzionali}
\label{sec:functional-requirements}

\begin{itemize}
    \item Invio e ricezione di messaggi tra nodi
    \item Scoperta automatica dei nodi
    \item Gestione della perdita di un nodo
\end{itemize}

\section{Requisiti non funzionali}
\label{sec:non-functional-requirements}

\begin{itemize}
    \item Portabilità tra Zenoh e Zenoh Pico
    \item Strategie alternative di gestione della perdita di nodi
\end{itemize}
